% Options for packages loaded elsewhere
\PassOptionsToPackage{unicode}{hyperref}
\PassOptionsToPackage{hyphens}{url}
\documentclass[
  12pt,
]{article}
\usepackage{xcolor}
\usepackage[margin=2.5cm]{geometry}
\usepackage{amsmath,amssymb}
\setcounter{secnumdepth}{5}
\usepackage{iftex}
\ifPDFTeX
  \usepackage[T1]{fontenc}
  \usepackage[utf8]{inputenc}
  \usepackage{textcomp} % provide euro and other symbols
\else % if luatex or xetex
  \usepackage{unicode-math} % this also loads fontspec
  \defaultfontfeatures{Scale=MatchLowercase}
  \defaultfontfeatures[\rmfamily]{Ligatures=TeX,Scale=1}
\fi
\usepackage{lmodern}
\ifPDFTeX\else
  % xetex/luatex font selection
\fi
% Use upquote if available, for straight quotes in verbatim environments
\IfFileExists{upquote.sty}{\usepackage{upquote}}{}
\IfFileExists{microtype.sty}{% use microtype if available
  \usepackage[]{microtype}
  \UseMicrotypeSet[protrusion]{basicmath} % disable protrusion for tt fonts
}{}
\makeatletter
\@ifundefined{KOMAClassName}{% if non-KOMA class
  \IfFileExists{parskip.sty}{%
    \usepackage{parskip}
  }{% else
    \setlength{\parindent}{0pt}
    \setlength{\parskip}{6pt plus 2pt minus 1pt}}
}{% if KOMA class
  \KOMAoptions{parskip=half}}
\makeatother
\usepackage{color}
\usepackage{fancyvrb}
\newcommand{\VerbBar}{|}
\newcommand{\VERB}{\Verb[commandchars=\\\{\}]}
\DefineVerbatimEnvironment{Highlighting}{Verbatim}{commandchars=\\\{\}}
% Add ',fontsize=\small' for more characters per line
\usepackage{framed}
\definecolor{shadecolor}{RGB}{248,248,248}
\newenvironment{Shaded}{\begin{snugshade}}{\end{snugshade}}
\newcommand{\AlertTok}[1]{\textcolor[rgb]{0.94,0.16,0.16}{#1}}
\newcommand{\AnnotationTok}[1]{\textcolor[rgb]{0.56,0.35,0.01}{\textbf{\textit{#1}}}}
\newcommand{\AttributeTok}[1]{\textcolor[rgb]{0.13,0.29,0.53}{#1}}
\newcommand{\BaseNTok}[1]{\textcolor[rgb]{0.00,0.00,0.81}{#1}}
\newcommand{\BuiltInTok}[1]{#1}
\newcommand{\CharTok}[1]{\textcolor[rgb]{0.31,0.60,0.02}{#1}}
\newcommand{\CommentTok}[1]{\textcolor[rgb]{0.56,0.35,0.01}{\textit{#1}}}
\newcommand{\CommentVarTok}[1]{\textcolor[rgb]{0.56,0.35,0.01}{\textbf{\textit{#1}}}}
\newcommand{\ConstantTok}[1]{\textcolor[rgb]{0.56,0.35,0.01}{#1}}
\newcommand{\ControlFlowTok}[1]{\textcolor[rgb]{0.13,0.29,0.53}{\textbf{#1}}}
\newcommand{\DataTypeTok}[1]{\textcolor[rgb]{0.13,0.29,0.53}{#1}}
\newcommand{\DecValTok}[1]{\textcolor[rgb]{0.00,0.00,0.81}{#1}}
\newcommand{\DocumentationTok}[1]{\textcolor[rgb]{0.56,0.35,0.01}{\textbf{\textit{#1}}}}
\newcommand{\ErrorTok}[1]{\textcolor[rgb]{0.64,0.00,0.00}{\textbf{#1}}}
\newcommand{\ExtensionTok}[1]{#1}
\newcommand{\FloatTok}[1]{\textcolor[rgb]{0.00,0.00,0.81}{#1}}
\newcommand{\FunctionTok}[1]{\textcolor[rgb]{0.13,0.29,0.53}{\textbf{#1}}}
\newcommand{\ImportTok}[1]{#1}
\newcommand{\InformationTok}[1]{\textcolor[rgb]{0.56,0.35,0.01}{\textbf{\textit{#1}}}}
\newcommand{\KeywordTok}[1]{\textcolor[rgb]{0.13,0.29,0.53}{\textbf{#1}}}
\newcommand{\NormalTok}[1]{#1}
\newcommand{\OperatorTok}[1]{\textcolor[rgb]{0.81,0.36,0.00}{\textbf{#1}}}
\newcommand{\OtherTok}[1]{\textcolor[rgb]{0.56,0.35,0.01}{#1}}
\newcommand{\PreprocessorTok}[1]{\textcolor[rgb]{0.56,0.35,0.01}{\textit{#1}}}
\newcommand{\RegionMarkerTok}[1]{#1}
\newcommand{\SpecialCharTok}[1]{\textcolor[rgb]{0.81,0.36,0.00}{\textbf{#1}}}
\newcommand{\SpecialStringTok}[1]{\textcolor[rgb]{0.31,0.60,0.02}{#1}}
\newcommand{\StringTok}[1]{\textcolor[rgb]{0.31,0.60,0.02}{#1}}
\newcommand{\VariableTok}[1]{\textcolor[rgb]{0.00,0.00,0.00}{#1}}
\newcommand{\VerbatimStringTok}[1]{\textcolor[rgb]{0.31,0.60,0.02}{#1}}
\newcommand{\WarningTok}[1]{\textcolor[rgb]{0.56,0.35,0.01}{\textbf{\textit{#1}}}}
\usepackage{graphicx}
\makeatletter
\newsavebox\pandoc@box
\newcommand*\pandocbounded[1]{% scales image to fit in text height/width
  \sbox\pandoc@box{#1}%
  \Gscale@div\@tempa{\textheight}{\dimexpr\ht\pandoc@box+\dp\pandoc@box\relax}%
  \Gscale@div\@tempb{\linewidth}{\wd\pandoc@box}%
  \ifdim\@tempb\p@<\@tempa\p@\let\@tempa\@tempb\fi% select the smaller of both
  \ifdim\@tempa\p@<\p@\scalebox{\@tempa}{\usebox\pandoc@box}%
  \else\usebox{\pandoc@box}%
  \fi%
}
% Set default figure placement to htbp
\def\fps@figure{htbp}
\makeatother
\setlength{\emergencystretch}{3em} % prevent overfull lines
\providecommand{\tightlist}{%
  \setlength{\itemsep}{0pt}\setlength{\parskip}{0pt}}
\usepackage{tabu}
\usepackage{longtable}
\usepackage{amsmath}
\usepackage{amssymb}
\usepackage{booktabs}
\usepackage{graphicx}
\usepackage{float}
\usepackage{setspace}
\usepackage{ragged2e}
\onehalfspacing
\justifying
\usepackage{booktabs}
\usepackage{longtable}
\usepackage{array}
\usepackage{multirow}
\usepackage{wrapfig}
\usepackage{float}
\usepackage{colortbl}
\usepackage{pdflscape}
\usepackage{tabu}
\usepackage{threeparttable}
\usepackage{threeparttablex}
\usepackage[normalem]{ulem}
\usepackage{makecell}
\usepackage{xcolor}
\usepackage{bookmark}
\IfFileExists{xurl.sty}{\usepackage{xurl}}{} % add URL line breaks if available
\urlstyle{same}
\hypersetup{
  pdftitle={Taller 3 -- plan de muestreo estándar},
  pdfauthor={Diego Fernando Muñoz Portela 2415620, Valery Rivera Lopez 2417038, Alejandro Barros Rayo 2415837},
  hidelinks,
  pdfcreator={LaTeX via pandoc}}

\title{Taller 3 -- plan de muestreo estándar}
\author{Diego Fernando Muñoz Portela 2415620, Valery Rivera Lopez
2417038, Alejandro Barros Rayo 2415837}
\date{14 de marzo de 2026}

\begin{document}
\maketitle

{
\setcounter{tocdepth}{3}
\tableofcontents
}
\section{Punto 1}\label{punto-1}

\subsection{Enunciado}\label{enunciado}

El jefe de compras de una empresa ha exigido un AQL de 2.5\% para la
característica de un cierto producto que se recibe en lotes de 3000
unidades y se desea obtener un plan de inspección por muestreo usando
las normas Militar Standard.

\subsection{Preguntas}\label{preguntas}

\subsubsection{a. ¿Cuál es el plan de muestreo simple adecuado bajo los
niveles de inspección I, II, III, en los niveles normal, reducido y
severo?. Evalué la conveniencia de los planes resultantes. (OC, AOQ,
ATI).}\label{a.-cuuxe1l-es-el-plan-de-muestreo-simple-adecuado-bajo-los-niveles-de-inspecciuxf3n-i-ii-iii-en-los-niveles-normal-reducido-y-severo.-evaluuxe9-la-conveniencia-de-los-planes-resultantes.-oc-aoq-ati.}

Primero buuscamos cada uno de los \(n\),\(c\) y \(r\) para cada uno de
los planes simples (1,2,3) en cada uno de los niveles por medio de la
funcion:

Representamos los n, r y c segun su clasificación de nivel en la
siguente tabla:

\begin{table}
\centering
\caption{\label{tab:tabla de c, n y r}Planes de muestreo simple según nivel de inspección}
\centering
\fontsize{12}{14}\selectfont
\begin{tabular}[t]{>{\centering\arraybackslash}p{2cm}>{\centering\arraybackslash}p{2cm}>{\centering\arraybackslash}p{2cm}>{\centering\arraybackslash}p{2cm}}
\toprule
\multicolumn{1}{c}{ } & \multicolumn{3}{c}{Parámetros del plan} \\
\cmidrule(l{3pt}r{3pt}){2-4}
Nivel & n & c & r\\
\midrule
\addlinespace[0.3em]
\multicolumn{4}{l}{\cellcolor[HTML]{f0f0f0}{\textbf{Normal}}}\\
\hspace{1em}I & 50 & 3 & 4\\
\hspace{1em}II & 125 & 7 & 8\\
\hspace{1em}III & 200 & 10 & 11\\
\addlinespace[0.3em]
\multicolumn{4}{l}{\cellcolor[HTML]{f0f0f0}{\textbf{Estricto}}}\\
\hspace{1em}I & 50 & 2 & 3\\
\hspace{1em}II & 125 & 5 & 6\\
\hspace{1em}III & 200 & 8 & 9\\
\addlinespace[0.3em]
\multicolumn{4}{l}{\cellcolor[HTML]{f0f0f0}{\textbf{Reducido}}}\\
\hspace{1em}I & 20 & 1 & 4\\
\hspace{1em}II & 50 & 3 & 6\\
\hspace{1em}III & 80 & 5 & 8\\
\bottomrule
\multicolumn{4}{l}{\rule{0pt}{1em}\textit{Note: }}\\
\multicolumn{4}{l}{\rule{0pt}{1em}n: tamaño de muestra; c: número de aceptación; r: número de rechazo}\\
\end{tabular}
\end{table}

\begin{center}\includegraphics{TALLER-3_files/figure-latex/Curvas OC 1-1} \end{center}

\paragraph{Análisis de las curvas
OC}\label{anuxe1lisis-de-las-curvas-oc}

La curva OC permite visualizar el comportamiento de los nueve planes de
muestreo simple obtenidos a partir de los parámetros mostrados en la
tabla para distintos niveles de inspección (estricto, normal y reducido)
y tres niveles de plan (I, II y III). En la gráfica se observa
claramente cómo la probabilidad de aceptación disminuye a medida que
aumenta la proporción defectuosa del lote, lo cual es consistente con la
lógica del muestreo de aceptación: cuanto peor es la calidad del lote,
menor es la probabilidad de que sea aceptado. Sin embargo, la rapidez
con la que disminuye esta probabilidad depende del tipo de inspección y
del tamaño de la muestra utilizado en cada plan.

Los planes estrictos (representados en rojo) muestran las pendientes más
pronunciadas dentro de la gráfica. Esto significa que, a medida que la
proporción de defectuosos aumenta ligeramente, la probabilidad de
aceptación cae rápidamente. Este comportamiento se explica porque los
valores de aceptación 𝑐 c son más pequeños que en los otros tipos de
inspección, lo que hace que el criterio para aceptar un lote sea más
exigente. En consecuencia, estos planes protegen principalmente al
consumidor, ya que reducen la posibilidad de aceptar lotes con niveles
de defectos relativamente altos. No obstante, esta mayor exigencia
también implica un mayor riesgo para el productor, puesto que incluso
pequeñas variaciones en la calidad del proceso pueden conducir al
rechazo del lote.

En contraste, los planes de inspección reducida (representados en verde)
presentan curvas más suaves y extendidas. En estos planes la
probabilidad de aceptación se mantiene relativamente alta incluso cuando
la proporción de defectos empieza a aumentar. Esto ocurre porque el
tamaño de muestra es menor y los números de aceptación son relativamente
más permisivos en comparación con los planes estrictos. Como resultado,
estos planes favorecen al productor, ya que disminuyen la probabilidad
de rechazo de sus lotes cuando la calidad del proceso es estable. Sin
embargo, esta ventaja implica que el consumidor puede estar más expuesto
a recibir lotes con niveles de defectos superiores al deseado, razón por
la cual este tipo de inspección generalmente se aplica únicamente cuando
el proveedor ha demostrado un desempeño de calidad consistente.

Por su parte, los planes de inspección normal (representados en azul) se
ubican en una posición intermedia entre los planes estrictos y los
reducidos. Sus curvas muestran un equilibrio entre la rapidez de rechazo
y la probabilidad de aceptación, lo que refleja el propósito de este
tipo de inspección como el esquema estándar utilizado en condiciones
normales de producción. Además, dentro de cada tipo de inspección se
observa que al pasar del plan I al plan III las curvas se vuelven más
empinadas. Esto se debe al incremento en el tamaño de la muestra (n),
que proporciona mayor información sobre la calidad del lote y aumenta la
capacidad del plan para distinguir entre lotes aceptables y no
aceptables.

En conjunto, la gráfica OC refleja claramente la lógica de los planes de
muestreo simple: los planes estrictos priorizan la protección del
consumidor, los planes reducidos favorecen la eficiencia para el
productor cuando existe confianza en la calidad del proceso, y los
planes normales mantienen un balance entre ambos intereses. Asimismo, el
aumento en el nivel del plan implica un mayor tamaño de muestra y, por
lo tanto, una mayor capacidad para discriminar entre diferentes niveles
de calidad en los lotes inspeccionados.

\begin{center}\includegraphics{TALLER-3_files/figure-latex/Curvas de salida promedio-1} \end{center}

\paragraph{Análisis de las curvas
AOQ}\label{anuxe1lisis-de-las-curvas-aoq}

La curva AOQ (Average Outgoing Quality) permite analizar la calidad
promedio de salida de los lotes después del proceso de inspección con
rectificación, considerando los nueve planes de muestreo simple
mostrados en la gráfica. En el eje horizontal se representa la
proporción defectuosa original del lote, mientras que en el eje vertical
se muestra la calidad promedio de salida (AOQ), es decir, el porcentaje
esperado de defectos que finalmente llega al cliente después de aplicar
el plan de muestreo. La forma característica de estas curvas muestra que
la calidad de salida primero aumenta, alcanza un máximo y luego
disminuye, debido a que cuando la proporción de defectos en el lote se
vuelve muy alta, los lotes tienden a ser rechazados con mayor frecuencia
y, al rectificarse, los defectos se eliminan antes de que el lote sea
enviado.

En la gráfica se observa que los planes de inspección estricta (rojos)
presentan valores máximos de AOQ más bajos y curvas que alcanzan su
punto máximo más rápidamente. Esto significa que la peor calidad
promedio que podría llegar al consumidor es relativamente pequeña, ya
que estos planes rechazan con mayor facilidad los lotes defectuosos.
Como consecuencia, los planes estrictos limitan el AOQ máximo (AOQL) y
brindan una mayor protección al consumidor. Sin embargo, esta misma
característica implica que los productores enfrentan una mayor
probabilidad de rechazo, lo que puede aumentar costos operativos debido
a inspecciones adicionales o reprocesamiento.

Por otro lado, los planes de inspección reducida (verdes) muestran los
valores más altos de AOQ y curvas más amplias. En estos casos, la
calidad promedio de salida puede alcanzar niveles más altos antes de
comenzar a disminuir, lo que indica que una mayor proporción de defectos
puede llegar al cliente antes de que los lotes sean rechazados
sistemáticamente. Esto ocurre porque los tamaños de muestra son menores
y el sistema de aceptación es más permisivo. En consecuencia, estos
planes favorecen al productor, ya que reducen la probabilidad de rechazo
del lote, pero implican un mayor riesgo potencial para el consumidor si
la calidad del proceso se deteriora.

Los planes de inspección normal (azules) nuevamente ocupan una posición
intermedia entre los planes estrictos y los reducidos. Sus curvas AOQ
alcanzan valores máximos moderados, reflejando un equilibrio entre la
protección al consumidor y la eficiencia para el productor. Este
comportamiento es consistente con el papel de la inspección normal como
el esquema estándar de control de calidad cuando no existen señales de
problemas significativos ni evidencia suficiente para aplicar inspección
estricta o reducida.

Asimismo, dentro de cada tipo de inspección se aprecia que al pasar del
plan I al plan III las curvas tienden a reducir ligeramente su AOQ
máximo y a concentrarse más cerca del origen. Esto se debe al aumento en
el tamaño de muestra, lo que mejora la capacidad del plan para detectar
lotes defectuosos y reducir la cantidad de defectos que finalmente
llegan al consumidor. Por tanto, los planes con mayor tamaño de muestra
ofrecen un mayor control sobre la calidad promedio de salida.

En conjunto, la gráfica AOQ complementa la interpretación de las curvas
OC, mostrando no solo la probabilidad de aceptar un lote, sino también
el nivel de calidad promedio que finalmente llega al cliente después del
proceso de inspección. Así, los planes estrictos minimizan la peor
calidad posible de salida, los planes reducidos privilegian la
eficiencia del productor cuando el proceso es confiable, y los planes
normales mantienen un equilibrio entre ambos enfoques dentro del sistema
de muestreo de aceptación.

\begin{center}\includegraphics{TALLER-3_files/figure-latex/inspeccion total-1} \end{center}

\paragraph{Análisis de las curvas
ATI}\label{anuxe1lisis-de-las-curvas-ati}

La curva ATI (Average Total Inspection) permite analizar el número
promedio de unidades inspeccionadas bajo los nueve planes de muestreo
simple considerados, teniendo en cuenta que cuando un lote es rechazado
se realiza inspección total con rectificación. En el eje horizontal se
representa la proporción defectuosa del lote, mientras que en el eje
vertical se muestra el número promedio de unidades inspeccionadas. La
gráfica evidencia que, a medida que aumenta la proporción de defectos,
el número promedio de unidades inspeccionadas también crece, hasta
aproximarse al tamaño total del lote. Esto ocurre porque los lotes con
mayor cantidad de defectos tienen una probabilidad más alta de ser
rechazados, lo que obliga a realizar una inspección completa para
eliminar las unidades defectuosas.

En la gráfica se observa que los planes de inspección estricta (rojos)
presentan un incremento más rápido en el ATI. Esto significa que incluso
para proporciones defectuosas relativamente pequeñas, estos planes
tienden a rechazar los lotes con mayor frecuencia, lo que conduce a
inspecciones más extensas. Como consecuencia, el número promedio de
unidades inspeccionadas se acerca rápidamente al tamaño total del lote.
Este comportamiento refleja nuevamente el carácter más exigente de la
inspección estricta, que busca proteger al consumidor al detectar con
mayor rapidez los lotes de calidad dudosa, aunque esto implique un mayor
esfuerzo de inspección para el productor.

Por el contrario, los planes de inspección reducida (verdes) muestran
curvas más graduales, lo que indica que el número promedio de unidades
inspeccionadas aumenta de forma más lenta conforme crece la proporción
defectuosa. Esto se debe a que estos planes aceptan los lotes con mayor
facilidad, por lo que en muchos casos solo se inspecciona la muestra
inicial y no se requiere realizar inspección completa del lote. En
consecuencia, estos planes reducen significativamente el esfuerzo de
inspección, lo que resulta ventajoso para el productor en términos de
costos y tiempo, aunque también implica una menor capacidad de detección
temprana de problemas de calidad.

Los planes de inspección normal (azules) nuevamente ocupan una posición
intermedia entre los estrictos y los reducidos. Sus curvas muestran un
crecimiento moderado del ATI, reflejando un equilibrio entre la
probabilidad de rechazo y el esfuerzo de inspección requerido. Este
comportamiento es coherente con el propósito de la inspección normal
como el esquema estándar utilizado cuando no existen evidencias claras
para aplicar reglas más severas o más permisivas.

Además, dentro de cada tipo de inspección se aprecia que al pasar del
plan I al plan III el ATI aumenta más rápidamente. Esto se debe al
incremento en el tamaño de muestra, lo que implica que incluso cuando el
lote es aceptado se inspecciona una mayor cantidad de unidades. Por lo
tanto, los planes con niveles más altos proporcionan una mayor
información sobre la calidad del lote, pero también requieren un mayor
esfuerzo de inspección.

En conjunto, la gráfica ATI complementa la interpretación de las curvas
OC y AOQ al mostrar el costo operativo del sistema de inspección en
términos del número de unidades inspeccionadas. Los planes estrictos
implican mayores niveles de inspección cuando la calidad del lote es
dudosa, los planes reducidos minimizan el esfuerzo de inspección cuando
existe confianza en el proceso productivo, y los planes normales ofrecen
un punto de equilibrio entre control de calidad y eficiencia operativa.

\textbackslash begin\{table\}{[}!h{]} \centering
\textbackslash caption\{\label{tab:tabla_comparativa}Tabla comparativa
de planes simples (AQL = 2.5\%, N = 3000)\} \centering
\resizebox{\ifdim\width>\linewidth\linewidth\else\width\fi}{!}{
\fontsize{10}{12}\selectfont
\begin{tabular}[t]{lccccccl}
\toprule
Inspeccion & Nivel & n & c & Pa (p=AQL) & AOQL & ATI (p=AQL) & Conveniencia\\
\midrule
\addlinespace[0.3em]
\multicolumn{8}{l}{\textbf{Normal}}\\
\hspace{1em}Normal & I & 50 & 3 & 96.4\% & 3.82\% & 157 & Balance costo-proteccion\\
\hspace{1em}Normal & II & 125 & 7 & 98.6\% & 3.45\% & 164 & Balance costo-proteccion\\
\hspace{1em}Normal & III & 200 & 10 & 98.7\% & 3.06\% & 235 & Balance costo-proteccion\\
\addlinespace[0.3em]
\multicolumn{8}{l}{\textbf{Estricta}}\\
\hspace{1em}Estricta & I & 50 & 2 & 87.1\% & 2.69\% & 432 & Alta proteccion, mayor costo\\
\hspace{1em}Estricta & II & 125 & 5 & 90.6\% & 2.44\% & 397 & Alta proteccion, mayor costo\\
\hspace{1em}Estricta & III & 200 & 8 & 93.4\% & 2.41\% & 384 & Alta proteccion, mayor costo\\
\addlinespace[0.3em]
\multicolumn{8}{l}{\textbf{Reducida}}\\
\hspace{1em}Reducida & I & 20 & 1 & 91.2\% & 4.11\% & 283 & Economico, historial confiable\\
\hspace{1em}Reducida & II & 50 & 3 & 96.4\% & 3.82\% & 157 & Economico, historial confiable\\
\hspace{1em}Reducida & III & 80 & 5 & 98.5\% & 3.87\% & 124 & Economico, historial confiable\\
\bottomrule
\end{tabular}} \textbackslash end\{table\}

\paragraph{Conveniencia de los planes}\label{conveniencia-de-los-planes}

Considerando el AQL de 2.5\% y lotes de N=3000:

\begin{enumerate}
\def\labelenumi{\arabic{enumi}.}
\item
  \textbf{Nivel I Normal} (n=50, c=3): conveniente cuando el costo de
  inspección es alto y el proveedor tiene historial confiable. Bajo
  costo pero menor discriminación.
\item
  \textbf{Nivel II Normal} (n=125, c=7): el plan de referencia según
  MIL-STD-105E; ofrece el mejor balance entre protección y costo de
  inspección.
\item
  \textbf{Nivel III Normal} (n=200, c=10): recomendable cuando las
  consecuencias de aceptar un lote defectuoso son graves. Mayor costo
  pero mayor protección.
\item
  \textbf{Planes Estrictos}: se activan cuando el proveedor ha tenido
  lotes rechazados consecutivos. Protegen mejor al consumidor pero
  incrementan el ATI.
\item
  \textbf{Planes Reducidos}: solo aplicables con historial comprobado de
  buena calidad. Reducen costos pero exponen al riesgo de mayor calidad
  de salida defectuosa (AOQL alto).
\end{enumerate}

\paragraph{Conclusión sobre los planes
simples}\label{conclusiuxf3n-sobre-los-planes-simples}

El plan de muestreo propuesto para el jefe de compras concluye que:El
Nivel II Normal es el punto de equilibrio óptimo: Ofrece una protección
estadística suficiente sin elevar los costos de inspección (\(n=125\)) a
niveles innecesarios para un inicio de relación comercial.El Sistema es
Preventivo, no solo Detectivo: La verdadera utilidad de estas curvas es
incentivar al proveedor a mantener la calidad. El aumento del costo de
inspección (ATI) y el endurecimiento de los criterios (Estricta) sirven
como penalización económica y técnica para el proveedor que no cumple
con el \(AQL\).Decisión Estratégica: Si el producto es crítico para la
seguridad o tiene un alto costo de falla, se debe optar por el Nivel
III, aceptando un mayor costo de inspección inicial a cambio de una
protección superior del consumidor

\subsubsection{b. Ubique ahora los respectivos planes de muestreo doble.
Evalúe la conveniencia de los planes resultantes. (OC, AOQ,
ATI).}\label{b.-ubique-ahora-los-respectivos-planes-de-muestreo-doble.-evaluxfae-la-conveniencia-de-los-planes-resultantes.-oc-aoq-ati.}

Realizamos:

Dado lo anterior, para este caso específico de inspección reducida,
trabajen con los siguientes planes identificados en las tablas (PDF)

\textbackslash begin\{table\}{[}!h{]} \centering
\textbackslash caption\{\label{tab:tabla_comparativa_doble}Planes de
muestreo doble (AQL = 2.5\%, N = 3000)\} \centering
\fontsize{10}{12}\selectfont

\begin{tabular}[t]{lccccccc}
\toprule
Inspeccion & Nivel & n1 & c1 & r1 & n2 & c2 & r2\\
\midrule
\addlinespace[0.3em]
\multicolumn{8}{l}{\textbf{Normal}}\\
\hspace{1em}Normal & I & 32 & 1 & 4 & 32 & 4 & 5\\
\hspace{1em}Normal & II & 80 & 3 & 7 & 80 & 8 & 9\\
\hspace{1em}Normal & III & 125 & 5 & 9 & 125 & 12 & 13\\
\addlinespace[0.3em]
\multicolumn{8}{l}{\textbf{Estricta}}\\
\hspace{1em}Estricta & I & 32 & 0 & 3 & 32 & 3 & 4\\
\hspace{1em}Estricta & II & 80 & 2 & 5 & 80 & 6 & 7\\
\hspace{1em}Estricta & III & 125 & 3 & 7 & 125 & 11 & 12\\
\addlinespace[0.3em]
\multicolumn{8}{l}{\textbf{Reducida}}\\
\hspace{1em}Reducida & I & 13 & 0 & 3 & 13 & 3 & 4\\
\hspace{1em}Reducida & II & 32 & 2 & 4 & 32 & 5 & 6\\
\hspace{1em}Reducida & III & 50 & 3 & 6 & 50 & 7 & 8\\
\bottomrule
\end{tabular}

\textbackslash end\{table\}

\begin{center}\includegraphics{TALLER-3_files/figure-latex/grafica OC-1} \end{center}

\paragraph{Análisis de las curvas
OC}\label{anuxe1lisis-de-las-curvas-oc-1}

La curva OC permite visualizar el comportamiento de los nueve planes de
muestreo doble obtenidos de las tablas MIL-STD-105E para un lote de 3000
unidades con AQL = 2.5\%. En ella se observa claramente cómo la
probabilidad de aceptación disminuye a medida que aumenta la proporción
de defectuosos en el lote, pero la velocidad de esa caída varía
notablemente según el tipo de inspección. Los planes estrictos (rojos)
son los que presentan una pendiente más pronunciada, lo que significa
que rechazan rápidamente los lotes con calidad dudosa. Esta
característica beneficia directamente al consumidor, ya que minimiza su
riesgo de aceptar lotes defectuosos (riesgo β), pero representa un
desafío para el productor, cuyos lotes tienen menos margen si su proceso
no es impecable.

En el extremo opuesto, los planes reducidos (verdes) muestran curvas
mucho más tendidas, manteniendo probabilidades de aceptación altas
incluso cuando la proporción de defectos supera el AQL pactado. Esto
favorece claramente al productor, reduciendo su riesgo de rechazo
injustificado (riesgo α), pero a costa de exponer al consumidor a una
mayor probabilidad de recibir lotes de mala calidad. Por esta razón, la
norma MIL-STD-105E solo permite la inspección reducida cuando el
proveedor ha demostrado un historial de calidad consistente y confiable
a lo largo del tiempo.

Los planes normales (azules) ocupan una posición intermedia,
equilibrando los intereses de ambas partes. A medida que se pasa del
nivel I al III, las curvas se vuelven más empinadas. Esto se debe al
aumento en el tamaño de muestra, lo que otorga mayor poder
discriminativo al plan y permite diferenciar con más claridad entre
lotes aceptables y rechazables. Así, la gráfica OC refleja fielmente la
filosofía del estándar, planes estrictos para proteger al consumidor
cuando hay dudas, reducidos para premiar al productor cuando hay
confianza, y normales como punto de partida equilibrado.

\begin{center}\includegraphics{TALLER-3_files/figure-latex/grafica AOQ-1} \end{center}

\paragraph{Análisis de las curvas
AOQ}\label{anuxe1lisis-de-las-curvas-aoq-1}

La curva AOQ muestra la calidad promedio que finalmente llega al
consumidor después de aplicar la inspección rectificadora, es decir,
asumiendo que los lotes rechazados se inspeccionan al 100\% y se
reemplazan las unidades defectuosas. Su forma característica de campana
invertida revela un punto máximo, conocido como AOQL, que representa la
peor calidad de salida que puede esperarse bajo ese plan,
independientemente de la calidad de entrada. En esta gráfica se observa
que los planes estrictos (rojos) alcanzan los valores más bajos de AOQL,
lo que garantiza una protección sólida al consumidor incluso en las
condiciones más desfavorables.

Los planes normales (azules) presentan AOQL moderados, ofreciendo un
equilibrio razonable entre la calidad de salida y el esfuerzo de
inspección. Son la opción más adecuada cuando no se tiene información
concluyente sobre la confiabilidad del proveedor, ya que aseguran que la
calidad promedio no se dispare sin incurrir en costos excesivos. Por su
parte, los planes reducidos (verdes), especialmente en el nivel I,
exhiben los AOQL más altos. Esto implica que, bajo este esquema, se
permite que una mayor proporción de defectuosos pueda escapar a la
inspección y llegar al consumidor.

Este comportamiento es coherente con la lógica del estándar, la
inspección reducida solo debe aplicarse cuando existe evidencia sólida
de que el proveedor entrega calidad superior de manera consistente,
porque relaja los controles justamente cuando el riesgo de encontrar
defectos es bajo. En cambio, cuando se sospecha que la calidad puede
haberse deteriorado, los planes estrictos actúan como un filtro más
riguroso para mantener la AOQ bajo control. La gráfica AOQ permite, por
tanto, cuantificar esa diferencia y elegir el plan más adecuado según el
nivel de exigencia deseado.

\paragraph{Conclusión sobre los planes
dobles}\label{conclusiuxf3n-sobre-los-planes-dobles}

En síntesis, los planes de muestreo doble analizados ofrecen muchas
opciones que se adaptan a diferentes escenarios y necesidades. Los
planes estrictos son la herramienta más eficaz para proteger al
consumidor cuando la confianza en el proveedor es limitada o cuando las
consecuencias de aceptar lotes defectuosos son graves, aunque esto
implique un mayor costo de inspección. Los planes reducidos, en cambio,
reducen drásticamente la carga de trabajo y los costos, pero solo son
recomendables cuando el proveedor ha demostrado un desempeño
sobresaliente y se puede asumir un mayor riesgo en la calidad de salida.

Los planes normales, y en particular el nivel II, representan el punto
de equilibrio clásico que la norma MIL-STD-105E propone como estándar.
Combinan una capacidad de discriminación aceptable con un esfuerzo de
inspección moderado, lo que los convierte en la opción más versátil y
utilizada en la práctica. En definitiva, la elección entre uno u otro
tipo de plan dependerá del balance que se quiera lograr entre la
protección al consumidor, los costos de inspección y la confianza
depositada en el proveedor, siendo las curvas OC y AOQ herramientas
fundamentales para tomar esa decisión con criterio técnico.

\subsection{c.~Comparación del mejor plan simple vs.~el mejor plan
doble}\label{c.-comparaciuxf3n-del-mejor-plan-simple-vs.-el-mejor-plan-doble}

Para realizar la comparación se selecciona el plan Normal Nivel II en
ambas familias, por ser el plan de referencia que establece la norma
MIL-STD-105E para condiciones estándar de operación. Los planes a
comparar son:

\begin{itemize}
  \item \textbf{Plan simple Normal II:} $n = 125$, $c = 7$, $r = 8$
  \item \textbf{Plan doble Normal II:} $n_1 = 80$, $c_1 = 3$, $r_1 = 7$; 
        $n_2 = 80$, $c_2 = 8$, $r_2 = 9$
\end{itemize}

\begin{Shaded}
\begin{Highlighting}[]
\FunctionTok{library}\NormalTok{(AcceptanceSampling)}
\FunctionTok{library}\NormalTok{(tidyverse)}

\NormalTok{N\_lote }\OtherTok{\textless{}{-}} \DecValTok{3000}
\NormalTok{pd     }\OtherTok{\textless{}{-}} \FunctionTok{seq}\NormalTok{(}\DecValTok{0}\NormalTok{, }\FloatTok{0.15}\NormalTok{, }\AttributeTok{length.out =} \DecValTok{200}\NormalTok{)}

\CommentTok{\# {-}{-}{-} Pa simple {-}{-}{-}}
\NormalTok{pa\_simple }\OtherTok{\textless{}{-}} \FunctionTok{OC2c}\NormalTok{(}\AttributeTok{n =} \DecValTok{125}\NormalTok{, }\AttributeTok{c =} \DecValTok{7}\NormalTok{, }\AttributeTok{type =} \StringTok{"binomial"}\NormalTok{, }\AttributeTok{pd =}\NormalTok{ pd)}\SpecialCharTok{@}\NormalTok{paccept}

\CommentTok{\# {-}{-}{-} Pa doble {-}{-}{-}}
\NormalTok{Pa\_doble }\OtherTok{\textless{}{-}} \ControlFlowTok{function}\NormalTok{(n1, c1, r1, n2, c2, p) \{}
\NormalTok{  pa1 }\OtherTok{\textless{}{-}} \FunctionTok{pbinom}\NormalTok{(c1, n1, p)}
\NormalTok{  pa2 }\OtherTok{\textless{}{-}} \FunctionTok{numeric}\NormalTok{(}\FunctionTok{length}\NormalTok{(p))}
  \ControlFlowTok{for}\NormalTok{ (i }\ControlFlowTok{in} \FunctionTok{seq\_along}\NormalTok{(p)) \{}
\NormalTok{    s }\OtherTok{\textless{}{-}} \DecValTok{0}
    \ControlFlowTok{for}\NormalTok{ (d1 }\ControlFlowTok{in}\NormalTok{ (c1 }\SpecialCharTok{+} \DecValTok{1}\NormalTok{)}\SpecialCharTok{:}\NormalTok{(r1 }\SpecialCharTok{{-}} \DecValTok{1}\NormalTok{)) \{}
      \ControlFlowTok{for}\NormalTok{ (d2 }\ControlFlowTok{in} \DecValTok{0}\SpecialCharTok{:}\NormalTok{(c2 }\SpecialCharTok{{-}}\NormalTok{ d1)) \{}
        \ControlFlowTok{if}\NormalTok{ (d2 }\SpecialCharTok{\textgreater{}=} \DecValTok{0} \SpecialCharTok{\&\&}\NormalTok{ d2 }\SpecialCharTok{\textless{}=}\NormalTok{ n2)}
\NormalTok{          s }\OtherTok{\textless{}{-}}\NormalTok{ s }\SpecialCharTok{+} \FunctionTok{dbinom}\NormalTok{(d1, n1, p[i]) }\SpecialCharTok{*} \FunctionTok{dbinom}\NormalTok{(d2, n2, p[i])}
\NormalTok{      \}}
\NormalTok{    \}}
\NormalTok{    pa2[i] }\OtherTok{\textless{}{-}}\NormalTok{ s}
\NormalTok{  \}}
\NormalTok{  pa1 }\SpecialCharTok{+}\NormalTok{ pa2}
\NormalTok{\}}
\NormalTok{pa\_doble }\OtherTok{\textless{}{-}} \FunctionTok{Pa\_doble}\NormalTok{(}\DecValTok{80}\NormalTok{, }\DecValTok{3}\NormalTok{, }\DecValTok{7}\NormalTok{, }\DecValTok{80}\NormalTok{, }\DecValTok{8}\NormalTok{, pd)}

\CommentTok{\# {-}{-}{-} ASN doble {-}{-}{-}}
\NormalTok{ASN\_doble }\OtherTok{\textless{}{-}} \ControlFlowTok{function}\NormalTok{(n1, c1, r1, n2, p) \{}
\NormalTok{  p\_segunda }\OtherTok{\textless{}{-}} \FunctionTok{numeric}\NormalTok{(}\FunctionTok{length}\NormalTok{(p))}
  \ControlFlowTok{for}\NormalTok{ (i }\ControlFlowTok{in} \FunctionTok{seq\_along}\NormalTok{(p)) \{}
\NormalTok{    s }\OtherTok{\textless{}{-}} \DecValTok{0}
    \ControlFlowTok{for}\NormalTok{ (d1 }\ControlFlowTok{in}\NormalTok{ (c1 }\SpecialCharTok{+} \DecValTok{1}\NormalTok{)}\SpecialCharTok{:}\NormalTok{(r1 }\SpecialCharTok{{-}} \DecValTok{1}\NormalTok{))}
\NormalTok{      s }\OtherTok{\textless{}{-}}\NormalTok{ s }\SpecialCharTok{+} \FunctionTok{dbinom}\NormalTok{(d1, n1, p[i])}
\NormalTok{    p\_segunda[i] }\OtherTok{\textless{}{-}}\NormalTok{ s}
\NormalTok{  \}}
\NormalTok{  n1 }\SpecialCharTok{+}\NormalTok{ p\_segunda }\SpecialCharTok{*}\NormalTok{ n2}
\NormalTok{\}}
\NormalTok{asn }\OtherTok{\textless{}{-}} \FunctionTok{ASN\_doble}\NormalTok{(}\DecValTok{80}\NormalTok{, }\DecValTok{3}\NormalTok{, }\DecValTok{7}\NormalTok{, }\DecValTok{80}\NormalTok{, pd)}

\CommentTok{\# {-}{-}{-} AOQ {-}{-}{-}}
\NormalTok{aoq\_simple }\OtherTok{\textless{}{-}}\NormalTok{ pa\_simple }\SpecialCharTok{*}\NormalTok{ pd }\SpecialCharTok{*}\NormalTok{ (N\_lote }\SpecialCharTok{{-}} \DecValTok{125}\NormalTok{) }\SpecialCharTok{/}\NormalTok{ N\_lote}
\NormalTok{aoq\_doble  }\OtherTok{\textless{}{-}}\NormalTok{ pa\_doble  }\SpecialCharTok{*}\NormalTok{ pd }\SpecialCharTok{*}\NormalTok{ (N\_lote }\SpecialCharTok{{-}}\NormalTok{ asn)  }\SpecialCharTok{/}\NormalTok{ N\_lote}

\CommentTok{\# {-}{-}{-} Data frames para ggplot {-}{-}{-}}
\NormalTok{df\_oc  }\OtherTok{\textless{}{-}} \FunctionTok{tibble}\NormalTok{(pd, }\AttributeTok{Simple =}\NormalTok{ pa\_simple, }\AttributeTok{Doble =}\NormalTok{ pa\_doble) }\SpecialCharTok{\%\textgreater{}\%}
            \FunctionTok{pivot\_longer}\NormalTok{(}\SpecialCharTok{{-}}\NormalTok{pd, }\AttributeTok{names\_to =} \StringTok{"Plan"}\NormalTok{, }\AttributeTok{values\_to =} \StringTok{"Pa"}\NormalTok{)}

\NormalTok{df\_aoq }\OtherTok{\textless{}{-}} \FunctionTok{tibble}\NormalTok{(pd, }\AttributeTok{Simple =}\NormalTok{ aoq\_simple, }\AttributeTok{Doble =}\NormalTok{ aoq\_doble) }\SpecialCharTok{\%\textgreater{}\%}
            \FunctionTok{pivot\_longer}\NormalTok{(}\SpecialCharTok{{-}}\NormalTok{pd, }\AttributeTok{names\_to =} \StringTok{"Plan"}\NormalTok{, }\AttributeTok{values\_to =} \StringTok{"AOQ"}\NormalTok{)}

\NormalTok{df\_asn }\OtherTok{\textless{}{-}} \FunctionTok{tibble}\NormalTok{(pd,}
                 \StringTok{\textasciigrave{}}\AttributeTok{Simple (n fijo)}\StringTok{\textasciigrave{}} \OtherTok{=} \FunctionTok{rep}\NormalTok{(}\DecValTok{125}\NormalTok{, }\FunctionTok{length}\NormalTok{(pd)),}
                 \StringTok{\textasciigrave{}}\AttributeTok{Doble (ASN)}\StringTok{\textasciigrave{}}     \OtherTok{=}\NormalTok{ asn) }\SpecialCharTok{\%\textgreater{}\%}
            \FunctionTok{pivot\_longer}\NormalTok{(}\SpecialCharTok{{-}}\NormalTok{pd, }\AttributeTok{names\_to =} \StringTok{"Plan"}\NormalTok{, }\AttributeTok{values\_to =} \StringTok{"Muestra"}\NormalTok{)}

\NormalTok{cols }\OtherTok{\textless{}{-}} \FunctionTok{c}\NormalTok{(}\StringTok{"Simple"} \OtherTok{=} \StringTok{"\#185FA5"}\NormalTok{, }\StringTok{"Doble"} \OtherTok{=} \StringTok{"\#3B6D11"}\NormalTok{,}
          \StringTok{"Simple (n fijo)"} \OtherTok{=} \StringTok{"\#185FA5"}\NormalTok{, }\StringTok{"Doble (ASN)"} \OtherTok{=} \StringTok{"\#3B6D11"}\NormalTok{)}
\end{Highlighting}
\end{Shaded}

\begin{Shaded}
\begin{Highlighting}[]
\FunctionTok{ggplot}\NormalTok{(df\_oc, }\FunctionTok{aes}\NormalTok{(}\AttributeTok{x =}\NormalTok{ pd, }\AttributeTok{y =}\NormalTok{ Pa, }\AttributeTok{color =}\NormalTok{ Plan, }\AttributeTok{linetype =}\NormalTok{ Plan)) }\SpecialCharTok{+}
  \FunctionTok{geom\_line}\NormalTok{(}\AttributeTok{size =} \FloatTok{1.1}\NormalTok{) }\SpecialCharTok{+}
  \FunctionTok{geom\_vline}\NormalTok{(}\AttributeTok{xintercept =} \FloatTok{0.025}\NormalTok{, }\AttributeTok{linetype =} \StringTok{"dotted"}\NormalTok{,}
             \AttributeTok{color =} \StringTok{"\#BA7517"}\NormalTok{, }\AttributeTok{size =} \FloatTok{0.8}\NormalTok{) }\SpecialCharTok{+}
  \FunctionTok{annotate}\NormalTok{(}\StringTok{"text"}\NormalTok{, }\AttributeTok{x =} \FloatTok{0.027}\NormalTok{, }\AttributeTok{y =} \FloatTok{0.3}\NormalTok{, }\AttributeTok{label =} \StringTok{"AQL = 2.5\%"}\NormalTok{,}
           \AttributeTok{color =} \StringTok{"\#BA7517"}\NormalTok{, }\AttributeTok{size =} \FloatTok{3.2}\NormalTok{, }\AttributeTok{hjust =} \DecValTok{0}\NormalTok{) }\SpecialCharTok{+}
  \FunctionTok{scale\_color\_manual}\NormalTok{(}\AttributeTok{values =}\NormalTok{ cols) }\SpecialCharTok{+}
  \FunctionTok{scale\_linetype\_manual}\NormalTok{(}\AttributeTok{values =} \FunctionTok{c}\NormalTok{(}\StringTok{"Simple"} \OtherTok{=} \StringTok{"solid"}\NormalTok{, }\StringTok{"Doble"} \OtherTok{=} \StringTok{"dashed"}\NormalTok{)) }\SpecialCharTok{+}
  \FunctionTok{scale\_x\_continuous}\NormalTok{(}\AttributeTok{labels =}\NormalTok{ scales}\SpecialCharTok{::}\NormalTok{percent) }\SpecialCharTok{+}
  \FunctionTok{scale\_y\_continuous}\NormalTok{(}\AttributeTok{labels =}\NormalTok{ scales}\SpecialCharTok{::}\NormalTok{percent, }\AttributeTok{limits =} \FunctionTok{c}\NormalTok{(}\DecValTok{0}\NormalTok{, }\DecValTok{1}\NormalTok{)) }\SpecialCharTok{+}
  \FunctionTok{labs}\NormalTok{(}\AttributeTok{title =} \StringTok{"Curva OC – Simple vs. Doble (Normal Nivel II)"}\NormalTok{,}
       \AttributeTok{x =} \StringTok{"Proporción defectuosa (p)"}\NormalTok{,}
       \AttributeTok{y =} \StringTok{"Probabilidad de aceptación"}\NormalTok{,}
       \AttributeTok{color =} \StringTok{"Plan"}\NormalTok{, }\AttributeTok{linetype =} \StringTok{"Plan"}\NormalTok{) }\SpecialCharTok{+}
  \FunctionTok{theme\_minimal}\NormalTok{() }\SpecialCharTok{+}
  \FunctionTok{theme}\NormalTok{(}\AttributeTok{plot.title =} \FunctionTok{element\_text}\NormalTok{(}\AttributeTok{hjust =} \FloatTok{0.5}\NormalTok{, }\AttributeTok{face =} \StringTok{"bold"}\NormalTok{, }\AttributeTok{color =} \StringTok{"brown"}\NormalTok{),}
        \AttributeTok{legend.position =} \StringTok{"bottom"}\NormalTok{)}
\end{Highlighting}
\end{Shaded}

\begin{center}\includegraphics{TALLER-3_files/figure-latex/grafica oc c-1} \end{center}

\begin{center}\rule{0.5\linewidth}{0.5pt}\end{center}

\textbf{Análisis OC} :

Las curvas OC de ambos planes son prácticamente equivalentes en la zona
del AQL (p = 2.5\%), donde ambos alcanzan una probabilidad de aceptación
superior al 98\%. Esto confirma que los dos planes ofrecen la misma
protección nominal al productor en condiciones de proceso controladas.
La diferencia entre ellos no radica en la capacidad de discriminación
---que es similar--- sino en cómo se distribuye el esfuerzo de muestreo
para lograr esa discriminación, aspecto que se revela con mayor claridad
en la curva ASN.

\begin{Shaded}
\begin{Highlighting}[]
\FunctionTok{ggplot}\NormalTok{(df\_aoq, }\FunctionTok{aes}\NormalTok{(}\AttributeTok{x =}\NormalTok{ pd, }\AttributeTok{y =}\NormalTok{ AOQ, }\AttributeTok{color =}\NormalTok{ Plan, }\AttributeTok{linetype =}\NormalTok{ Plan)) }\SpecialCharTok{+}
  \FunctionTok{geom\_line}\NormalTok{(}\AttributeTok{size =} \FloatTok{1.1}\NormalTok{) }\SpecialCharTok{+}
  \FunctionTok{geom\_vline}\NormalTok{(}\AttributeTok{xintercept =} \FloatTok{0.025}\NormalTok{, }\AttributeTok{linetype =} \StringTok{"dotted"}\NormalTok{,}
             \AttributeTok{color =} \StringTok{"\#BA7517"}\NormalTok{, }\AttributeTok{size =} \FloatTok{0.8}\NormalTok{) }\SpecialCharTok{+}
  \FunctionTok{scale\_color\_manual}\NormalTok{(}\AttributeTok{values =}\NormalTok{ cols) }\SpecialCharTok{+}
  \FunctionTok{scale\_linetype\_manual}\NormalTok{(}\AttributeTok{values =} \FunctionTok{c}\NormalTok{(}\StringTok{"Simple"} \OtherTok{=} \StringTok{"solid"}\NormalTok{, }\StringTok{"Doble"} \OtherTok{=} \StringTok{"dashed"}\NormalTok{)) }\SpecialCharTok{+}
  \FunctionTok{scale\_x\_continuous}\NormalTok{(}\AttributeTok{labels =}\NormalTok{ scales}\SpecialCharTok{::}\NormalTok{percent) }\SpecialCharTok{+}
  \FunctionTok{scale\_y\_continuous}\NormalTok{(}\AttributeTok{labels =}\NormalTok{ scales}\SpecialCharTok{::}\NormalTok{percent) }\SpecialCharTok{+}
  \FunctionTok{labs}\NormalTok{(}\AttributeTok{title =} \StringTok{"Curva AOQ – Simple vs. Doble (Normal Nivel II)"}\NormalTok{,}
       \AttributeTok{x =} \StringTok{"Proporción defectuosa (p)"}\NormalTok{,}
       \AttributeTok{y =} \StringTok{"Calidad de salida promedio (AOQ)"}\NormalTok{,}
       \AttributeTok{color =} \StringTok{"Plan"}\NormalTok{, }\AttributeTok{linetype =} \StringTok{"Plan"}\NormalTok{) }\SpecialCharTok{+}
  \FunctionTok{theme\_minimal}\NormalTok{() }\SpecialCharTok{+}
  \FunctionTok{theme}\NormalTok{(}\AttributeTok{plot.title =} \FunctionTok{element\_text}\NormalTok{(}\AttributeTok{hjust =} \FloatTok{0.5}\NormalTok{, }\AttributeTok{face =} \StringTok{"bold"}\NormalTok{, }\AttributeTok{color =} \StringTok{"brown"}\NormalTok{),}
        \AttributeTok{legend.position =} \StringTok{"bottom"}\NormalTok{)}
\end{Highlighting}
\end{Shaded}

\begin{center}\includegraphics{TALLER-3_files/figure-latex/grafica aoq c-1} \end{center}

\begin{center}\rule{0.5\linewidth}{0.5pt}\end{center}

\textbf{Análisis AOQ} :

Ambos planes presentan un AOQL similar, ubicado alrededor de p =
0.04--0.05. El plan simple tiene un AOQL ligeramente más bajo porque su
tamaño de muestra fijo de 125 unidades garantiza una inspección
constante independientemente del resultado. El plan doble, al poder
tomar decisiones con solo \(n_1 = 80\) unidades cuando el lote es
claramente bueno o malo, inspecciona menos en promedio y permite que
marginalmente más defectuosos escapen sin ser detectados. En la práctica
esta diferencia es pequeña y operativamente insignificante para el rango
de calidad esperado cerca del AQL.

\begin{Shaded}
\begin{Highlighting}[]
\FunctionTok{ggplot}\NormalTok{(df\_asn, }\FunctionTok{aes}\NormalTok{(}\AttributeTok{x =}\NormalTok{ pd, }\AttributeTok{y =}\NormalTok{ Muestra, }\AttributeTok{color =}\NormalTok{ Plan, }\AttributeTok{linetype =}\NormalTok{ Plan)) }\SpecialCharTok{+}
  \FunctionTok{geom\_line}\NormalTok{(}\AttributeTok{size =} \FloatTok{1.1}\NormalTok{) }\SpecialCharTok{+}
  \FunctionTok{geom\_vline}\NormalTok{(}\AttributeTok{xintercept =} \FloatTok{0.025}\NormalTok{, }\AttributeTok{linetype =} \StringTok{"dotted"}\NormalTok{,}
             \AttributeTok{color =} \StringTok{"\#BA7517"}\NormalTok{, }\AttributeTok{size =} \FloatTok{0.8}\NormalTok{) }\SpecialCharTok{+}
  \FunctionTok{scale\_color\_manual}\NormalTok{(}\AttributeTok{values =}\NormalTok{ cols) }\SpecialCharTok{+}
  \FunctionTok{scale\_linetype\_manual}\NormalTok{(}\AttributeTok{values =} \FunctionTok{c}\NormalTok{(}\StringTok{"Simple (n fijo)"} \OtherTok{=} \StringTok{"solid"}\NormalTok{,}
                                   \StringTok{"Doble (ASN)"}     \OtherTok{=} \StringTok{"dashed"}\NormalTok{)) }\SpecialCharTok{+}
  \FunctionTok{scale\_x\_continuous}\NormalTok{(}\AttributeTok{labels =}\NormalTok{ scales}\SpecialCharTok{::}\NormalTok{percent) }\SpecialCharTok{+}
  \FunctionTok{labs}\NormalTok{(}\AttributeTok{title =} \StringTok{"Tamaño de muestra promedio (ASN) – Simple vs. Doble"}\NormalTok{,}
       \AttributeTok{x =} \StringTok{"Proporción defectuosa (p)"}\NormalTok{,}
       \AttributeTok{y =} \StringTok{"Unidades inspeccionadas en promedio"}\NormalTok{,}
       \AttributeTok{color =} \StringTok{"Plan"}\NormalTok{, }\AttributeTok{linetype =} \StringTok{"Plan"}\NormalTok{) }\SpecialCharTok{+}
  \FunctionTok{theme\_minimal}\NormalTok{() }\SpecialCharTok{+}
  \FunctionTok{theme}\NormalTok{(}\AttributeTok{plot.title =} \FunctionTok{element\_text}\NormalTok{(}\AttributeTok{hjust =} \FloatTok{0.5}\NormalTok{, }\AttributeTok{face =} \StringTok{"bold"}\NormalTok{, }\AttributeTok{color =} \StringTok{"brown"}\NormalTok{),}
        \AttributeTok{legend.position =} \StringTok{"bottom"}\NormalTok{)}
\end{Highlighting}
\end{Shaded}

\begin{center}\includegraphics{TALLER-3_files/figure-latex/grafica asn-1} \end{center}

\begin{center}\rule{0.5\linewidth}{0.5pt}\end{center}

\textbf{Análisis ASN} :

Esta gráfica es la clave de la comparación. El plan simple inspecciona
siempre \(n = 125\) unidades, sin importar qué tan bueno o malo sea el
lote. El plan doble, en cambio, tiene un tamaño de muestra esperado
(ASN) variable:

\begin{itemize}
  \item Cuando $p$ es baja (lote claramente bueno), la primera muestra de 
        $n_1 = 80$ unidades suele ser suficiente para aceptar, y el ASN se 
        aproxima a 80.
  \item Cuando $p$ es alta (lote claramente malo), la primera muestra también 
        basta para rechazar, y el ASN nuevamente se acerca a 80.
  \item En la zona de incertidumbre (alrededor de $p \approx 0.04$--$0.06$), 
        se requiere la segunda muestra con frecuencia, y el ASN puede llegar 
        hasta 160 unidades.
\end{itemize}

Esto implica que el plan doble es más eficiente que el simple cuando el
proceso opera en los extremos (muy buena o muy mala calidad), pero puede
requerir más inspección que el simple en la zona intermedia. Para un
proveedor con proceso estable cerca del AQL, el ASN del doble ronda las
90--100 unidades, lo que representa un ahorro de aproximadamente 20\%
frente al plan simple.

\begin{Shaded}
\begin{Highlighting}[]
\FunctionTok{library}\NormalTok{(knitr)}
\FunctionTok{library}\NormalTok{(kableExtra)}

\CommentTok{\# Precalcular valores antes del data.frame}
\NormalTok{pa\_aql\_simple }\OtherTok{\textless{}{-}} \FunctionTok{round}\NormalTok{(}\FunctionTok{OC2c}\NormalTok{(}\DecValTok{125}\NormalTok{, }\DecValTok{7}\NormalTok{, }\AttributeTok{type =} \StringTok{"binomial"}\NormalTok{, }
                             \AttributeTok{pd =} \FloatTok{0.025}\NormalTok{)}\SpecialCharTok{@}\NormalTok{paccept }\SpecialCharTok{*} \DecValTok{100}\NormalTok{, }\DecValTok{1}\NormalTok{)}
\NormalTok{pa\_aql\_doble  }\OtherTok{\textless{}{-}} \FunctionTok{round}\NormalTok{(}\FunctionTok{Pa\_doble}\NormalTok{(}\DecValTok{80}\NormalTok{, }\DecValTok{3}\NormalTok{, }\DecValTok{7}\NormalTok{, }\DecValTok{80}\NormalTok{, }\DecValTok{8}\NormalTok{, }\FloatTok{0.025}\NormalTok{) }\SpecialCharTok{*} \DecValTok{100}\NormalTok{, }\DecValTok{1}\NormalTok{)}
\NormalTok{aoql\_simple   }\OtherTok{\textless{}{-}} \FunctionTok{round}\NormalTok{(}\FunctionTok{max}\NormalTok{(aoq\_simple) }\SpecialCharTok{*} \DecValTok{100}\NormalTok{, }\DecValTok{2}\NormalTok{)}
\NormalTok{aoql\_doble    }\OtherTok{\textless{}{-}} \FunctionTok{round}\NormalTok{(}\FunctionTok{max}\NormalTok{(aoq\_doble) }\SpecialCharTok{*} \DecValTok{100}\NormalTok{, }\DecValTok{2}\NormalTok{)}
\NormalTok{asn\_aql       }\OtherTok{\textless{}{-}} \FunctionTok{round}\NormalTok{(}\FunctionTok{ASN\_doble}\NormalTok{(}\DecValTok{80}\NormalTok{, }\DecValTok{3}\NormalTok{, }\DecValTok{7}\NormalTok{, }\DecValTok{80}\NormalTok{, }\FloatTok{0.025}\NormalTok{), }\DecValTok{0}\NormalTok{)}

\NormalTok{tabla\_c }\OtherTok{\textless{}{-}} \FunctionTok{data.frame}\NormalTok{(}
  \AttributeTok{Criterio =} \FunctionTok{c}\NormalTok{(}\StringTok{"Tamaño de muestra"}\NormalTok{, }
               \StringTok{"Pa en AQL (p=2.5\%)"}\NormalTok{,}
               \StringTok{"AOQL"}\NormalTok{, }
               \StringTok{"ASN en AQL"}\NormalTok{, }
               \StringTok{"Complejidad operativa"}\NormalTok{,}
               \StringTok{"Psicología del proveedor"}\NormalTok{,}
               \StringTok{"Recomendado cuando..."}\NormalTok{),}
  \AttributeTok{Simple =} \FunctionTok{c}\NormalTok{(}\StringTok{"Fijo: n=125"}\NormalTok{,}
             \FunctionTok{paste0}\NormalTok{(pa\_aql\_simple, }\StringTok{"\%"}\NormalTok{),}
             \FunctionTok{paste0}\NormalTok{(aoql\_simple, }\StringTok{"\%"}\NormalTok{),}
             \StringTok{"125 (siempre)"}\NormalTok{,}
             \StringTok{"Baja — una sola decisión"}\NormalTok{,}
             \StringTok{"Neutral"}\NormalTok{,}
             \StringTok{"Proceso variable o errático"}\NormalTok{),}
  \AttributeTok{Doble  =} \FunctionTok{c}\NormalTok{(}\StringTok{"Variable: 80–160"}\NormalTok{,}
             \FunctionTok{paste0}\NormalTok{(pa\_aql\_doble, }\StringTok{"\%"}\NormalTok{),}
             \FunctionTok{paste0}\NormalTok{(aoql\_doble, }\StringTok{"\%"}\NormalTok{),}
             \FunctionTok{paste0}\NormalTok{(}\StringTok{"\textasciitilde{}"}\NormalTok{, asn\_aql),}
             \StringTok{"Media — dos etapas posibles"}\NormalTok{,}
             \StringTok{"Favorable (segunda oportunidad)"}\NormalTok{,}
             \StringTok{"Proceso estable, costo inspección alto"}\NormalTok{)}
\NormalTok{)}

\FunctionTok{kbl}\NormalTok{(tabla\_c,}
    \AttributeTok{caption   =} \StringTok{"Comparación del mejor plan simple vs. el mejor plan doble (Normal Nivel II)"}\NormalTok{,}
    \AttributeTok{col.names =} \FunctionTok{c}\NormalTok{(}\StringTok{"Criterio"}\NormalTok{, }\StringTok{"Simple (n=125, c=7)"}\NormalTok{, }\StringTok{"Doble (n1=80, c1=3)"}\NormalTok{),}
    \AttributeTok{align     =} \FunctionTok{c}\NormalTok{(}\StringTok{"l"}\NormalTok{,}\StringTok{"c"}\NormalTok{,}\StringTok{"c"}\NormalTok{),}
    \AttributeTok{booktabs  =} \ConstantTok{TRUE}\NormalTok{) }\SpecialCharTok{\%\textgreater{}\%}
  \FunctionTok{kable\_styling}\NormalTok{(}\AttributeTok{font\_size =} \DecValTok{10}\NormalTok{,}
                \AttributeTok{latex\_options =} \FunctionTok{c}\NormalTok{(}\StringTok{"hold\_position"}\NormalTok{, }\StringTok{"scale\_down"}\NormalTok{))}
\end{Highlighting}
\end{Shaded}

\begin{table}[!h]
\centering
\caption{\label{tab:tabla comparativa c}Comparación del mejor plan simple vs. el mejor plan doble (Normal Nivel II)}
\centering
\resizebox{\ifdim\width>\linewidth\linewidth\else\width\fi}{!}{
\fontsize{10}{12}\selectfont
\begin{tabular}[t]{lcc}
\toprule
Criterio & Simple (n=125, c=7) & Doble (n1=80, c1=3)\\
\midrule
Tamaño de muestra & Fijo: n=125 & Variable: 80–160\\
Pa en AQL (p=2.5\%) & 98.6\% & 98.3\%\\
AOQL & 3.45\% & 3.25\%\\
ASN en AQL & 125 (siempre) & \textasciitilde{}91\\
Complejidad operativa & Baja — una sola decisión & Media — dos etapas posibles\\
\addlinespace
Psicología del proveedor & Neutral & Favorable (segunda oportunidad)\\
Recomendado cuando... & Proceso variable o errático & Proceso estable, costo inspección alto\\
\bottomrule
\end{tabular}}
\end{table}

\begin{center}\rule{0.5\linewidth}{0.5pt}\end{center}

\textbf{Conclusión final del punto c:} En síntesis, ambos planes ofrecen
una protección equivalente al consumidor en términos de curva OC y AOQL.
La diferencia determinante es el esfuerzo de inspección: el plan doble
permite reducir el tamaño de muestra esperado cuando el proceso del
proveedor es estable, a cambio de una mayor complejidad administrativa.

El \textbf{plan simple Normal II} es preferible cuando el proceso del
proveedor es errático o cuando la simplicidad operativa es prioritaria,
ya que su tamaño de muestra fijo garantiza decisiones consistentes sin
importar el estado del lote.

El \textbf{plan doble Normal II} es preferible cuando el costo unitario
de inspección es significativo y el proveedor tiene un historial de
calidad demostrado, pues su ASN promedio cercano al AQL puede ser hasta
un 20\% menor que el del plan simple, traduciendo ese ahorro
directamente en reducción de costos operativos sin sacrificar la
protección exigida.

\end{document}
